\documentclass[11pt,a4paper]{article}
\usepackage[T1]{fontenc}
\usepackage[utf8]{inputenc}
\usepackage{lmodern,microtype}
\usepackage[margin=25mm,headheight=15pt]{geometry}
\usepackage{amsmath,amssymb,amsthm,mathtools}
\usepackage{booktabs,array,longtable}
\usepackage{xcolor,enumitem,listings,fancyhdr,needspace}
\usepackage[hyphens]{url}
\usepackage{hyperref}
\definecolor{inkblue}{HTML}{173B57}
\definecolor{softblue}{HTML}{EEF4F8}
\hypersetup{colorlinks=true,linkcolor=inkblue,citecolor=inkblue,urlcolor=inkblue,
 pdftitle={Critical Values and Reciprocal Discriminants},
 pdfauthor={Research manuscript prepared for Keio}}
\pagestyle{fancy}
\fancyhf{}
\fancyhead[L]{\small\textsc{Critical values and reciprocal discriminants}}
\fancyhead[R]{\small 23 November 2026}
\fancyfoot[C]{\thepage}
\renewcommand{\headrulewidth}{0.3pt}
\setlength{\parskip}{4pt}
\setlength{\parindent}{1.25em}
\setlength{\emergencystretch}{1.5em}
\setlist[enumerate]{itemsep=3pt,topsep=4pt}
\setlist[itemize]{itemsep=3pt,topsep=4pt}
\newtheorem{theorem}{Theorem}[section]
\newtheorem{lemma}[theorem]{Lemma}
\newtheorem{proposition}[theorem]{Proposition}
\newtheorem{corollary}[theorem]{Corollary}
\theoremstyle{definition}
\newtheorem{definition}[theorem]{Definition}
\theoremstyle{remark}
\newtheorem{remark}[theorem]{Remark}
\newcommand{\QQ}{\mathbb Q}
\newcommand{\ZZ}{\mathbb Z}
\newcommand{\CC}{\mathbb C}
\newcommand{\FF}{\mathbb F}
\newcommand{\PP}{\mathbb P}
\newcommand{\Disc}{\operatorname{Disc}}
\newcommand{\Res}{\operatorname{Res}}
\newcommand{\den}{\operatorname{den}}
\newcommand{\ord}{\operatorname{ord}}
\newcommand{\lc}{\operatorname{lc}}
\newcommand{\status}[1]{\par\medskip\noindent\fcolorbox{inkblue}{softblue}{%
 \parbox{\dimexpr\linewidth-2\fboxsep-2\fboxrule\relax}{#1}}\par\medskip}
\lstset{basicstyle=\ttfamily\footnotesize,breaklines=true,columns=fullflexible,
 keepspaces=true,showstringspaces=false,frame=single,rulecolor=\color{inkblue},
 xleftmargin=0pt,xrightmargin=0pt,tabsize=4}

\title{\textbf{Critical Values and Reciprocal Discriminants}\\[6pt]
 \large A proof of Jones's Conjecture 3.3 in its integral-parameter setting\\[3pt]
 \normalsize With a global integral-weight extension of a quadrinomial gcd law}
\author{\normalsize Research manuscript prepared for Keio}
\date{23 November 2026}

\begin{document}
\maketitle
\thispagestyle{empty}

\begin{abstract}
Let $q\ge5$ be prime, $d=(q-1)/2$, and let $H_q\in\ZZ[u]$ be the monic
polynomial determined by $x^dH_q(x+x^{-1})=\Phi_q(x)$. We prove that
\[
 q^{1-d}\Disc_u\bigl(H_q(u)+q(Y-1)\bigr)
\]
is an Eisenstein polynomial at $q$. Consequently, the normalized
discriminants in Jones's 2021 Conjecture~3.3 are primitive irreducible
integer polynomials of degree $(q-3)/2$ for every prime $q\ge5$ and
every nonzero integer parameter $r$. Together with a self-contained
proof of the reciprocal discriminant formula, this proves the full
conjecture in the integral-parameter formulation used in the subsequent
literature. The factorization formula itself was proved previously by
Barman, Narode and Wagh. Our proof of irreducibility passes from the
exact valuation $1/(d-1)$ of a shifted critical point to the exact
valuation $d/(d-1)$ of its shifted critical value. An independent proof
recovers every critical point rationally from the square of its critical
value. As a second application of exact valuation and trace arguments,
we prove the quadrinomial gcd law for every integer weight other than
$0$ and $1$, and obtain a denominator obstruction for rational weights.
All assertions are proved throughout the stated parameter ranges; finite exact
computations serve only as checks. Historical priority is not certified
by the literature search reported here.
\end{abstract}

\status{\textbf{Scope.} This manuscript contains complete mathematical
arguments, not a finite-height reduction or a conditional proof.
The claims concern the stated polynomial families. The separate
existence problem for infinitely many squarefree prime specializations
is not resolved. This is a research manuscript, not an externally
refereed publication.}

\begingroup
\small\setcounter{tocdepth}{1}\setlength{\parskip}{0pt}
\makeatletter
\renewcommand*\l@section{\@dottedtocline{1}{0em}{1.5em}}
\makeatother
\tableofcontents
\endgroup
\newpage

\section{The existing problem and the exact conclusions}\label{sec:target}

Jones \cite[Conjecture~3.3, p.~1476]{Jones} considered, for an odd prime
$q$, $a\in\{0,1\}$, and a nonzero integer $r$, the reciprocal family
\begin{equation}\label{eq:jones-family}
 \mathcal F_{2^a q,t}(x)=\Phi_{2^a q}(x)+4rq^2t\,x^{(q-1)/2}.
\end{equation}
Here $t$ is an indeterminate and $\Phi_N$ is the $N$th cyclotomic
polynomial. His conjecture specifies the exact discriminant
factorization and the degree, primitivity, and irreducibility of its
normalized half-degree discriminant factors.

The 2026 paper of Barman, Narode and Wagh
\cite[Theorem~1.3]{BNW} proves the discriminant factorization and
explicitly describes its contribution as a partial proof. The
irreducibility assertion is distinct from that formula. We prove it
below for every odd prime and include the formula's proof to make the
completion self-contained. The related monogeneity criterion of
Narode \cite{Narode} concerns a different conclusion and is not used
as an irreducibility theorem for these factors.

\paragraph{Parameter convention.}
Jones writes ``arbitrary nonzero constant'' for $r$ after introducing
integer coefficients. We use the explicit convention
$r\in\ZZ\setminus\{0\}$ in \cite[(1.4)--(1.6)]{BNW}.
Irreducibility extends to every $r\in\QQ^\times$, but integer
coefficients need not: at $q=5,r=1/25$ the normalized factor is
$1-16t/5$. Thus the integer-parameter hypothesis is stated rather
than silently extended.

\paragraph{Conventions.}
For a monic degree-$n$ polynomial with roots $\alpha_1,\ldots,\alpha_n$,
\[
 \Disc(f)=\prod_{i<j}(\alpha_i-\alpha_j)^2
 =(-1)^{n(n-1)/2}\Res(f,f').
\]
The discriminant of a monic linear polynomial is $1$.
The content of an integer polynomial is the positive gcd of its
coefficients. Irreducibility is over $\QQ$, unless another field is
specified. The symbol $\doteq$ means equality up to a nonzero rational
constant. We normalize $v_p(p)=1$ on algebraic extensions of $\QQ_p$.

Define the monic trace polynomials by
\begin{equation}\label{eq:trace-recurrence}
 C_0(u)=2,\qquad C_1(u)=u,\qquad
 C_{j+1}(u)=uC_j(u)-C_{j-1}(u).
\end{equation}
For $j\ge1$, $C_j$ is monic of degree $j$, and
$C_j(z+z^{-1})=z^j+z^{-j}$. Put
\begin{equation}\label{eq:H-definition}
 d=\frac{q-1}{2},\qquad \varepsilon=(-1)^d,\qquad
 H_q(u)=1+\sum_{j=1}^{d}C_j(u),\qquad
 D_q(Z)=\Disc_u(H_q(u)+Z).
\end{equation}
The half-degree polynomials corresponding to
\eqref{eq:jones-family} are
\begin{equation}\label{eq:g-definition}
 g_0(u;t)=H_q(u)+4rq^2t,\qquad
 g_1(u;t)=\varepsilon H_q(-u)+4rq^2t.
\end{equation}
Both are monic of degree $d$ in $u$.

\begin{theorem}[Jones's discriminant conjecture]\label{thm:jones}
For every odd prime $q$, every $r\in\ZZ\setminus\{0\}$, and
$a\in\{0,1\}$, the polynomials in
\eqref{eq:jones-family}--\eqref{eq:g-definition} satisfy
\begin{align}
 \Disc_x(\mathcal F_{q,t})
 &=q(4qrt+1)(4q^2rt+\varepsilon)\Disc_u(g_0)^2,
 \label{eq:jones-zero}\\
 \Disc_x(\mathcal F_{2q,t})
 &=q(4qrt+\varepsilon)(4q^2rt+1)\Disc_u(g_1)^2.
 \label{eq:jones-one}
\end{align}
There exist $h_0,h_1\in\ZZ[t]$ such that
\begin{equation}\label{eq:normalized-h}
 \Disc_u(g_a)=q^{d-1}h_a(t),\qquad h_a(0)=1.
\end{equation}
For $q=3$, $h_0=h_1=1$. For $q\ge5$, each $h_a$ is primitive,
irreducible over $\QQ_q$ and over $\QQ$, and has degree $d-1$.
Moreover $h_1(t)=h_0(\varepsilon t)$.
\end{theorem}

The mechanism proving the nontrivial irreducibility assertion is the
stronger certificate
\begin{equation}\label{eq:E-certificate}
 \boxed{E_q(Y)=q^{1-d}D_q(q(Y-1))\in\ZZ[Y]
 \quad\text{is Eisenstein at }q.}
\end{equation}
Indeed $h_0(t)=E_q(1+4rqt)$. This affine substitution is invertible
over $\QQ$ and $\QQ_q$ because $r\ne0$; it preserves irreducibility.

\begin{theorem}[Integral weights]\label{thm:integer-main}
Let
\[
 A_0(t)=0,\quad A_1(t)=1,\quad
 A_{j+2}(t)=-A_{j+1}(t)-tA_j(t).
\]
For all integers $k>m\ge1$ and
$\lambda\in\ZZ\setminus\{0,1\}$, put $\delta=\gcd(k,m)$ and
\[
 F_\lambda(t)=A_k(t)^2-\lambda A_{k-m}(t)A_{k+m}(t).
\]
Then
\begin{equation}\label{eq:integral-main}
 \gcd_{\QQ[t]}(F_\lambda,F'_\lambda)\doteq A_\delta.
\end{equation}
In particular, $F_\lambda=A_\delta^2B_\lambda$ with
$B_\lambda\in\ZZ[t]$ squarefree and coprime to $A_\delta$,
$\deg B_\lambda=k-\delta$, and the smooth projective normalization of
$y^2=F_\lambda(t)$ has genus
$\lfloor(k-\delta-1)/2\rfloor$.
\end{theorem}

The weight $4$ case is the corrected gcd conjecture in \cite{KSL};
it was established in the preceding research manuscript \cite{Previous}.
The unrestricted integer-weight conclusion is a genuine enlargement
of that parameter range. We do not claim that the enlarged statement
was itself a previously named conjecture. Its proof is included in
Section~\ref{sec:integral-weights}.

\section{A critical-value form of Eisenstein's criterion}\label{sec:criterion}

The following lemma isolates the valuation calculation. It uses
classical nonarchimedean valuation properties; no new status is
claimed for Eisenstein's criterion itself.

\begin{theorem}[Critical-value certificate]\label{thm:criterion}
Let $p$ be prime and let
\[
 P(X)=X^n+\sum_{j=0}^{n-1}a_jX^j\in\ZZ[X],\qquad n\ge2.
\]
Assume $p\nmid n(n-1)$, $p\mid a_j$ for $1\le j\le n-1$, and
$v_p(a_1)=1$. There is no condition on $a_0$. Then
\begin{equation}\label{eq:general-certificate}
 \mathcal E_P(Y)=p^{-(n-1)}
 \Disc_X\bigl(P(X)-P(0)+pY\bigr)
\end{equation}
belongs to $\ZZ[Y]$, has degree $n-1$, and is Eisenstein at $p$.
Its leading coefficient is $(-1)^{n(n-1)/2}n^n$.
Consequently it is irreducible over $\QQ_p$ and $\QQ$.
\end{theorem}

\begin{proof}
Let $c$ be any root of $P'$ in $\overline{\QQ}_p$ and write $v=v_p$.
Since $a_1\ne0$, $c\ne0$. If $v(c)\le0$, the term $nc^{n-1}$
in $P'(c)$ has strictly smaller valuation than each other nonzero
term: for $1\le j\le n-1$,
\[
 v(ja_jc^{j-1})-v(nc^{n-1})
 \ge 1-(n-j)v(c)>0.
\]
A sum with a unique least-valuation term is nonzero, a contradiction.
Thus $v(c)>0$. The constant term $a_1$ has valuation $1$, whereas
every intermediate derivative term, with $j\ge2$, has valuation
at least $1+(j-1)v(c)>1$. Cancellation therefore requires
\begin{equation}\label{eq:critical-point-v}
 v(c)=\frac1{n-1}.
\end{equation}

Subtract $cP'(c)=0$ from $P(c)-a_0$ to obtain
\begin{equation}\label{eq:critical-value-identity}
 P(c)-a_0=-(n-1)c^n
 -\sum_{j=2}^{n-1}(j-1)a_jc^j.
\end{equation}
The first term has valuation $n/(n-1)$, using $p\nmid n-1$.
Each remaining nonzero term has valuation at least
$1+j/(n-1)>n/(n-1)$. Thus there is again a unique term of least
valuation, and this time it gives the exact equality
\begin{equation}\label{eq:critical-value-v}
 \boxed{v(P(c)-a_0)=\frac n{n-1}.}
\end{equation}

Let $c_1,\ldots,c_{n-1}$ be the roots of $P'$, counted with
multiplicity. The resultant formula gives
\begin{align*}
 \Disc_X(P(X)-a_0+pY)
 &=(-1)^{n(n-1)/2}n^n
   \prod_{i=1}^{n-1}(P(c_i)-a_0+pY).
\end{align*}
The sign from interchanging the two resultant arguments is $1$,
since $n(n-1)$ is even. Hence $\mathcal E_P$ has leading
coefficient $L=(-1)^{n(n-1)/2}n^n$, a $p$-adic unit, and roots
\[
 y_i=-\frac{P(c_i)-a_0}{p},\qquad v(y_i)=\frac1{n-1}.
\]
Write $\mathcal E_P(Y)=LY^{n-1}+b_1Y^{n-2}+\cdots+b_{n-1}$.
Initially $b_j\in\ZZ[1/p]$. The elementary symmetric expressions
in the $y_i$ show that $v(b_j)\ge j/(n-1)>0$ for $j\ge1$.
Since $b_j$ is rational, its valuation is an integer (or $+\infty$).
It follows that $b_j\in p\ZZ$: there were no possible denominator
primes other than $p$. For the constant term there is no sum,
so $v(b_{n-1})=1$ exactly. This proves the asserted Eisenstein
conditions and integrality.

For completeness, local irreducibility also follows directly from
these root valuations. A monic factor in $\QQ_p[Y]$ of positive
degree $s<n-1$ would have constant term of valuation $s/(n-1)$,
which is not an integer. This is impossible for a nonzero element
of $\QQ_p$. Thus the polynomial is irreducible over $\QQ_p$ and
therefore over $\QQ$.
\end{proof}

\begin{remark}[Content and the base field]
The general polynomial $\mathcal E_P$ need not be primitive over
$\ZZ$. For example, $n=2$, $p=3$, and $P(X)=X^2+6X$ give
$\mathcal E_P(Y)=12-4Y$. The theorem asserts irreducibility over
$\QQ_p$ and $\QQ$, where nonzero constants are units. The
primitivity needed for Jones's factors will follow separately from
their constant coefficient $1$.
\end{remark}

\begin{corollary}[A critical value generates the critical-point field]
\label{cor:critical-field}
Under the hypotheses of Theorem~\ref{thm:criterion}, if $P'(c)=0$,
then
\[
 \QQ_p(P(c))=\QQ_p(c),\qquad [\QQ_p(c):\QQ_p]=n-1.
\]
The corresponding equality $\QQ(P(c))=\QQ(c)$ holds for every
algebraic critical point in characteristic zero. Distinct critical
points have distinct critical values.
\end{corollary}

\begin{proof}
The derivative $P'$ has degree $n-1$, leading coefficient $n$,
and is Eisenstein at $p$, so its local and rational degrees are
$n-1$. The polynomial in \eqref{eq:general-certificate} is
irreducible of the same degree; its roots are affine transforms of
the critical values. Field inclusion in each direction of the
degree comparison gives the equalities. In characteristic zero an
irreducible polynomial is separable, so its $n-1$ displayed roots
are distinct. This also proves distinctness of critical values.
\end{proof}

\section{The cyclotomic trace polynomial}\label{sec:cyclotomic}

\begin{lemma}[Shift and first coefficient]\label{lem:shift}
For every odd prime $q$, the polynomial $H_q$ in
\eqref{eq:H-definition} satisfies
\begin{align}
 x^dH_q(x+x^{-1})&=\Phi_q(x),\label{eq:cyclotomic-trace}\\
 H_q(u)&\equiv(u-2)^d\pmod q,\label{eq:shift-congruence}\\
 H_q(2)&=q,\qquad H_q(-2)=(-1)^d,\label{eq:endpoints}\\
 H'_q(2)&=\frac{q(q^2-1)}{24}.\label{eq:first-coefficient}
\end{align}
For $q\ge5$, $v_q(H'_q(2))=1$.
\end{lemma}

\begin{proof}
The trace recurrence gives
\[
 x^dH_q(x+x^{-1})
 =x^d+\sum_{j=1}^d(x^{d+j}+x^{d-j})
 =1+x+\cdots+x^{2d}=\Phi_q(x).
\]
Modulo $q$ this is $(x-1)^{q-1}=(x-1)^{2d}$, which also equals
$x^d(x+x^{-1}-2)^d$. The substitution
$u\mapsto x+x^{-1}$ is injective from $\FF_q[u]$ into
$\FF_q[x,x^{-1}]$: the highest positive Laurent exponent of a
nonzero degree-$j$ polynomial is $j$, with the same nonzero
leading coefficient. This proves \eqref{eq:shift-congruence}.

By recurrence $C_j(2)=2$ and $C_j(-2)=2(-1)^j$, giving the
endpoint values. Differentiating the recurrence at $2$ shows by
induction that $C'_j(2)=j^2$. Consequently
\[
 H'_q(2)=\sum_{j=1}^dj^2
 =\frac{d(d+1)(2d+1)}6
 =\frac{q(q^2-1)}{24}.
\]
For $q\ge5$ neither $24$ nor $q^2-1$ is divisible by $q$.
\end{proof}

\begin{theorem}[An Eisenstein discriminant for every odd prime]
\label{thm:Eq}
For every prime $q\ge5$, the polynomial $E_q$ in
\eqref{eq:E-certificate} is Eisenstein at $q$ and has degree
$d-1$ and leading coefficient $(-1)^{d(d-1)/2}d^d$.
In particular, $D_q$ is irreducible over $\QQ_q$ and $\QQ$.
\end{theorem}

\begin{proof}
Set $P(X)=H_q(X+2)$. Lemma~\ref{lem:shift} gives
$P(X)\equiv X^d\pmod q$, $P(0)=q$, and $v_q(P'(0))=1$.
Since $2\le d<q$ and $1\le d-1<q$, $q\nmid d(d-1)$.
Theorem~\ref{thm:criterion} therefore applies with $p=q,n=d$.
Translation of $u$ does not change a discriminant, and
\[
 P(X)-P(0)+qY=H_q(X+2)+q(Y-1).
\]
Thus $\mathcal E_P(Y)=E_q(Y)$. The inverse affine substitution
and multiplication by $q^{d-1}$ recover $D_q$ from $E_q$.
\end{proof}

\paragraph{Where noncancellation enters.}
The congruence alone proves only the known Eisenstein behavior of
the shifted trace polynomial, as in \cite[Lemma~3.2]{Jones}.
What is needed for the conjecture is the exact additional calculation
\[
 P'(c)=0\quad\Longrightarrow\quad
 v_q(c)=\frac1{d-1},\qquad
 v_q(P(c)-q)=\frac d{d-1}.
\]
The second equality does not follow merely by substituting into
$P$ and estimating each term: the derivative equation must first
remove the linear term as in \eqref{eq:critical-value-identity}.
It is this equality that forces the full degree of the critical
value and establishes \eqref{eq:E-certificate}.

\section{Discriminants, primitivity, and completion of the conjecture}
\label{sec:completion}

\begin{lemma}[Reciprocal discriminant identity]\label{lem:reciprocal-disc}
If $g(u)$ is monic of degree $d$ and $f(x)=x^dg(x+x^{-1})$, then
\begin{equation}\label{eq:reciprocal-disc}
 \Disc_x(f)=g(2)g(-2)\Disc_u(g)^2.
\end{equation}
\end{lemma}

\begin{proof}
If the roots of $g$ are $\beta_i$, then
$f(x)=\prod_i f_i(x)$ with $f_i(x)=x^2-\beta_i x+1$.
Here $\Disc(f_i)=\beta_i^2-4$. Moreover, if $\rho$ is a root
of $f_i$, then $f_j(\rho)=(\beta_i-\beta_j)\rho$, and the
product of the two roots of $f_i$ is $1$. Thus
\[
 \Res(f_i,f_j)=(\beta_i-\beta_j)^2.
\]
The product formula for discriminants gives
\[
 \Disc(f)=\prod_i(\beta_i^2-4)
 \prod_{i<j}(\beta_i-\beta_j)^4.
\]
Its first product is $g(2)g(-2)$ and its second is $\Disc(g)^2$.
The identity holds initially for distinct roots and then identically
in the coefficients, including cases with repeated roots.
\end{proof}

\begin{lemma}[The constant coefficient]\label{lem:constant}
For every odd prime $q$, $\Disc(H_q)=q^{d-1}$.
\end{lemma}

\begin{proof}
Write $\zeta$ for a primitive $q$th root. The $d$ roots of $H_q$
are $\zeta^j+\zeta^{-j}=2\cos(2\pi j/q)$, $1\le j\le d$.
They are distinct and real, so $\Disc(H_q)>0$.

For a root $\alpha\ne1$ of $\Phi_q$, differentiating
$(x-1)\Phi_q(x)=x^q-1$ gives
$\Phi'_q(\alpha)=q\alpha^{q-1}/(\alpha-1)$.
The product of the roots $\alpha$ is $1$, and
$\prod_\alpha(\alpha-1)=\Phi_q(1)=q$, since $q-1$ is even.
It follows that $\prod_\alpha\Phi'_q(\alpha)=q^{q-2}$, and hence
\[
 \Disc(\Phi_q)=(-1)^dq^{2d-1}.
\]
Apply Lemma~\ref{lem:reciprocal-disc} to
\eqref{eq:cyclotomic-trace} and use \eqref{eq:endpoints}:
\[
 (-1)^dq^{2d-1}=q(-1)^d\Disc(H_q)^2.
\]
Taking the positive square root proves the claim, including $d=1$.
\end{proof}

\begin{proof}[Proof of Theorem~\ref{thm:jones}]
First suppose $q\ge5$. By Theorem~\ref{thm:Eq},
\[
 h_0(t):=q^{1-d}D_q(4rq^2t)=E_q(1+4rqt)
\]
is an integer polynomial of degree $d-1$, irreducible over
$\QQ_q$ and $\QQ$. Lemma~\ref{lem:constant} gives
$h_0(0)=q^{1-d}D_q(0)=1$. Its content is therefore $1$.

Since $g_1(u;t)=\varepsilon\bigl(H_q(-u)+\varepsilon4rq^2t\bigr)$,
reflection of the variable and multiplication by $\varepsilon$
do not change the discriminant: the powers in their transformation
laws are $d(d-1)$ and $2d-2$, both even. Thus
\[
 \Disc_u(g_1)=D_q(\varepsilon4rq^2t),\qquad
 h_1(t)=h_0(\varepsilon t).
\]
This proves all normalized-factor assertions for both $a$.

For $q=3$ the polynomials $g_0,g_1$ are monic linear, their
discriminants are $1$, and $h_0=h_1=1$.

It remains to establish the two exact factorizations, for all $q$.
Equation~\eqref{eq:cyclotomic-trace}, together with
$\Phi_{2q}(x)=\Phi_q(-x)$, gives
\[
 \mathcal F_{2^a q,t}(x)=x^dg_a(x+x^{-1};t).
\]
The endpoint values are
\begin{align*}
 g_0(2;t)&=q(1+4rqt),&g_0(-2;t)&=\varepsilon+4rq^2t,\\
 g_1(2;t)&=1+4rq^2t,&g_1(-2;t)&=q(\varepsilon+4rqt).
\end{align*}
Substitution in Lemma~\ref{lem:reciprocal-disc} gives
\eqref{eq:jones-zero} and \eqref{eq:jones-one} exactly.
\end{proof}

\begin{remark}[Coefficient normalization]
Our normalization is fixed by the standard discriminant and
$h_a(0)=1$. Some displayed factors in the computational tables
of \cite{Jones} have the opposite overall sign. Multiplying such
a factor by $-1$ has no effect on the squared factor in the reciprocal
discriminant formula, on content, or on irreducibility. The
definition \eqref{eq:normalized-h} fixes the sign throughout this
manuscript.
\end{remark}

\subsection{Small examples of the global certificate}

For $q=5$,
\[
 H_5(u)=u^2+u-1,\quad D_5(Z)=5-4Z,\quad E_5(Y)=5-4Y,
\]
and $h_0(t)=h_1(t)=1-80rt$.
For $q=7$,
\begin{align*}
 H_7(u)&=u^3+u^2-2u-1,\\
 D_7(Z)&=49+14Z-27Z^2,\\
 E_7(Y)&=-27Y^2+56Y-28,\\
 h_0(t)&=1+56rt-21168r^2t^2,\qquad h_1(t)=h_0(-t).
\end{align*}
The coefficients visibly satisfy the respective Eisenstein
conditions. These examples illustrate the universal argument;
the proof does not infer the all-prime assertion from them.

\section{A second proof by rational reconstruction}\label{sec:inverse}

\subsection{An explicit inverse from a critical value}
\label{subsec:critical-value-inverse}

The field generated by a critical value can also be identified directly,
without using its valuation.  The identity below supplies an independent
route from irreducibility of the derivative to irreducibility of the
discriminant polynomial.

\begin{lemma}[Recovery of a critical point]
\label{lem:critical-value-inverse}
Let $q\geq 3$ be odd, put $d=(q-1)/2$, and define
\[
 H_q(u)=1+\sum_{j=1}^{d}C_j(u),
 \qquad
 A_q(u)=2(q^2+1)-(q^2-1)u.
\]
The following identity holds in $\ZZ[u]$:
\begin{equation}
\label{eq:critical-value-polynomial-identity}
 \begin{split}
 A_q(u)H_q(u)^2
 &-4(4-u^2)H_q(u)H_q'(u)\\
 &+4(4-u^2)(2-u)H_q'(u)^2
 =4q^2.
 \end{split}
\end{equation}
Consequently, if $c\in\CC$ satisfies $H_q'(c)=0$, then $H_q(c)\neq 0$,
\begin{equation}
\label{eq:critical-value-rational-inverse}
 c=\frac{2(q^2+1)}{q^2-1}
   -\frac{4q^2}{(q^2-1)H_q(c)^2},
\end{equation}
and
\begin{equation}
\label{eq:critical-value-field-equality}
 \QQ(c)=\QQ\bigl(H_q(c)^2\bigr)=\QQ\bigl(H_q(c)\bigr).
\end{equation}
\end{lemma}

\begin{proof}
For $0<x<\pi/2$, set $u=2\cos(2x)$.  The finite geometric-series identity
for the Dirichlet kernel gives
\[
 H_q(u)=1+2\sum_{j=1}^{d}\cos(2jx)
       =\frac{\sin(qx)}{\sin x}.
\]
Differentiating $\sin(qx)=\sin x\,H_q(u)$ and using
$du/dx=-8\sin x\cos x$ yields
\[
 q\cos(qx)
 =\cos x\bigl(H_q(u)-2(2-u)H_q'(u)\bigr).
\]
Squaring this equation and the equation for $\sin(qx)$, and adding after
multiplying the latter square by $q^2$, gives
\[
 (2+u)\bigl(H_q(u)-2(2-u)H_q'(u)\bigr)^2
 +q^2(2-u)H_q(u)^2=4q^2.
\]
Expanding the left-hand side proves
\eqref{eq:critical-value-polynomial-identity} for every real $u\in(-2,2)$.
Both sides are polynomials, so the identity holds identically, and in
particular at every complex $u$.

At a critical point $c$, the identity reduces to
\[
 \bigl(2(q^2+1)-(q^2-1)c\bigr)H_q(c)^2=4q^2.
\]
Its right-hand side is nonzero.  Thus $H_q(c)\neq 0$, and solving for $c$
proves \eqref{eq:critical-value-rational-inverse}.  Since $H_q$ has rational
coefficients, this inverse gives the inclusions
\[
 \QQ(c)\subseteq\QQ\bigl(H_q(c)^2\bigr)
 \subseteq\QQ\bigl(H_q(c)\bigr)\subseteq\QQ(c),
\]
which prove \eqref{eq:critical-value-field-equality}.
\end{proof}

\begin{corollary}[A second proof of discriminant irreducibility]
\label{cor:critical-inverse-discriminant}
If $q\geq 5$ is prime, then
\[
 D_q(Z)=\Disc_u\bigl(H_q(u)+Z\bigr)
\]
is irreducible over $\QQ$ and has degree $d-1=(q-3)/2$.
\end{corollary}

\begin{proof}
Recall the shifted Eisenstein calculation:
\[
 H_q(X+2)\equiv X^d\pmod q,
 \qquad
 H_q'(2)=\sum_{j=1}^{d}j^2=\frac{q(q^2-1)}{24}.
\]
All nonleading coefficients of $H_q'(X+2)$ are therefore divisible by $q$;
its constant coefficient is divisible by $q$ but not by $q^2$, and its
leading coefficient $d$ is not divisible by $q$.  Hence $H_q'(X+2)$ is
$q$-Eisenstein.  It follows that $H_q'$ is irreducible over $\QQ$ and that
every root $c$ of $H_q'$ satisfies
\[
 [\QQ(c):\QQ]=d-1.
\]
Lemma~\ref{lem:critical-value-inverse} now gives
$[\QQ(H_q(c)):\QQ]=d-1$.

Writing $c_1,\ldots,c_{d-1}$ for the roots of $H_q'$, the resultant formula
for the discriminant gives
\[
 D_q(Z)=(-1)^{d(d-1)/2}d^d
        \prod_{i=1}^{d-1}\bigl(Z+H_q(c_i)\bigr).
\]
Thus $D_q$ has degree $d-1$ and vanishes at $-H_q(c)$.  The latter element
has degree $d-1$ over $\QQ$, so its minimal polynomial accounts for all of
$D_q$, up to a nonzero rational constant.  This proves irreducibility.
\end{proof}

\begin{remark}[Ordered absolute critical values]
\label{rem:absolute-critical-values}
For odd $q\geq5$, the roots of $H_q$ are the $d$ distinct real numbers
$2\cos(2\pi j/q)$, $1\leq j\leq d$.  Rolle's theorem therefore places all
$d-1$ roots of $H_q'$ in $(-2,2)$.  At each such root $c$, the identity above
reads
\[
 |H_q(c)|^2
 =\frac{4q^2}{2(q^2+1)-(q^2-1)c}.
\]
The right-hand side is strictly increasing as a function of $c\in(-2,2)$.
Consequently, distinct critical points have distinct absolute critical
values.  Equivalently, writing $c=2\cos(2x)$ with $0<x<\pi/2$, the absolute
value at a critical point is
\[
 \left|\frac{\sin(qx)}{\sin x}\right|
 =\frac{q}{\sqrt{1+(q^2-1)\sin^2x}},
\]
whose displayed right-hand side is strictly decreasing in $x$.  These
formulas apply at critical points; they are not formulas for $|H_q|$ at
arbitrary points.  For $q=3$, the derivative is constant and there are no
critical points; the discriminant polynomial is $D_3(Z)=1$.
\end{remark}


\section{Consequences and the precise remaining boundary}\label{sec:consequences}

\subsection{Local degree and full generic symmetric group}

The certificate proves more than rational irreducibility. Every
critical value of $H_q$ has degree $d-1$ over $\QQ_q$ and generates
the same local field as the associated critical point. This is the
field-degree obstruction in Corollary~\ref{cor:critical-field}.

\begin{corollary}[Generic Galois group]\label{cor:generic-galois}
For every odd prime $q$, the Galois group of $H_q(u)+T$ over
$\CC(T)$ is $S_d$, and the Galois group over $\QQ(T)$ is also $S_d$.
\end{corollary}

\begin{proof}
The case $d=1$ is immediate. For $d\ge2$, the derivative is separable
because $H'_q(X+2)$ is irreducible in characteristic zero. The
discriminant polynomial $D_q(T)$ is likewise separable by
Theorem~\ref{thm:Eq}. Therefore $H_q$ has $d-1$ simple critical
points with pairwise distinct critical values.

Consider the degree-$d$ map of complex projective lines
$u\mapsto-H_q(u)$. Removing its branch values and their preimages
gives a connected covering. Its monodromy group is transitive.
A small loop about each finite branch value exchanges exactly the
two sheets meeting at its unique simple critical point; it is a
transposition. These finite loops generate the monodromy, since a
loop about infinity is the inverse of their product. A transitive
permutation group generated by transpositions is $S_d$: the graph
whose edges are those transpositions is connected, and edge
transpositions of a connected graph generate all permutations.
This proves the geometric assertion. The arithmetic group contains
the geometric group and is a subgroup of $S_d$, so it too is $S_d$.
\end{proof}

This is a standard monodromy consequence of simple critical points
and distinct critical values. It is a generic assertion with
indeterminate $T$, not a claim that every rational specialization
has Galois group $S_d$. It is not presented as an additional
previously named conjecture.

\subsection{What changes in the monogeneity applications}

Theorem~\ref{thm:jones} removes the explicit assumption of
Conjecture~3.3 from \cite[Theorem~3.4]{Jones}. That theorem still
assumes the existence of appropriate squarefree values of its
auxiliary polynomials. Its deductions at such values are now
available without the discriminant conjecture, but the existence
and infinitude of suitable prime values remain separate questions.
In particular, this manuscript supplies no unconditional proof of
the number-field $abc$ conjecture or of infinitely many monogenic
prime specializations in every higher degree. The distinction is
also present in \cite[Theorem~1.4]{BNW}.

\section{Integral weights and rational denominator obstructions}
\label{sec:integral-weights}

The notation in this section is independent of the cyclotomic notation:
here $d$ denotes $\gcd(k,m)$, not $(q-1)/2$.

% A self-contained section-body fragment for incorporation into the main manuscript.
% Requires amsmath, amssymb, amsthm, and the theorem, lemma, proposition,
% corollary, and remark environments. All labels have the prefix iw:.
% Insert after the main manuscript's existing integral-weights section heading.
% No external mathematical statement from the main manuscript is assumed.

This section establishes a generalization of the weight-four identity.
The generalization is stated on its own mathematical merits; we do not
assert that the entire weighted family was previously posed as a named
conjecture, or that the results below have certified historical priority.

Define the integer polynomials
\begin{equation}\label{iw:A}
 A_0(t)=0,\qquad A_1(t)=1,\qquad
 A_{j+2}(t)=-A_{j+1}(t)-tA_j(t)\quad(j\ge0).
\end{equation}
For integers \(k>m\ge1\) and a rational number \(\lambda\), put
\begin{equation}\label{iw:F}
 F_{\lambda;k,m}(t)
   =A_k(t)^2-\lambda A_{k-m}(t)A_{k+m}(t).
\end{equation}
We will abbreviate this polynomial to \(F_\lambda\), and write
\begin{equation}\label{iw:parameters}
 d=\gcd(k,m),\qquad K=k/d,\qquad M=m/d,\qquad
 \Delta=K^2-M^2.
\end{equation}
Thus \(K>M\ge1\), \(\gcd(K,M)=1\), and \(\Delta\ge3\).
Association in \(\mathbb Q[t]\) is denoted by \(\doteq\).
Squarefreeness throughout this section means absence of repeated
complex roots, or equivalently squarefreeness over \(\mathbb Q\);
nonzero constant polynomials are squarefree.

\begin{theorem}[All integral weights]\label{iw:integral}
For every integer \(\lambda\notin\{0,1\}\) and every pair \(k>m\ge1\),
\begin{equation}\label{iw:integral-gcd}
 \boxed{\gcd_{\mathbb Q[t]}(F_\lambda,F_\lambda')\doteq A_d.}
\end{equation}
Every root of \(A_d\) is a double root of \(F_\lambda\), and every
other root of \(F_\lambda\) is simple. Moreover,
\begin{equation}\label{iw:integral-quotient}
 F_\lambda=A_d^2B_\lambda,\qquad B_\lambda\in\mathbb Z[t],
 \qquad \gcd(B_\lambda,B_\lambda')=
 \gcd(A_d,B_\lambda)=1,
\end{equation}
where the last two gcds are monic gcds over \(\mathbb Q\).
One has
\begin{equation}\label{iw:integral-degree-genus}
 \deg B_\lambda=k-d\ge1,\qquad
 g\bigl(\operatorname{Norm}(y^2=F_\lambda(t))\bigr)
 =\left\lfloor\frac{k-d-1}{2}\right\rfloor.
\end{equation}
Here \(\operatorname{Norm}\) means the smooth projective normalization,
and \(g\) is its geometric genus.
\end{theorem}

The assertion includes all negative integral weights, all prime
weights, and all integral square weights greater than one.
The proof will give a stronger necessary condition for an extra
multiple root when the coefficient is rational.

\subsection{The reciprocal representation and its exceptional points}

For later use, define
\begin{align}
 \phi(z)&=\frac{z}{(1+z)^2},\label{iw:phi}\\
 D_\lambda(z)
   &=(z^k-1)^2-\lambda(z^{k-m}-1)(z^{k+m}-1)\notag\\
   &=(1-\lambda)(z^{2k}+1)
        +\lambda(z^{k+m}+z^{k-m})-2z^k,\label{iw:D}\\
 H_\lambda(z)
   &=z^{-k}D_\lambda(z)\notag\\
   &=(1-\lambda)(z^k+z^{-k})
       +\lambda(z^m+z^{-m})-2.\label{iw:H}
\end{align}

\begin{lemma}[Exact substitution]\label{iw:substitution}
For \(j\ge1\) and \(z\ne\pm1\),
\begin{equation}\label{iw:A-substitution}
 A_j(\phi(z))
 =(-1)^{j-1}\frac{z^j-1}{(z-1)(z+1)^{j-1}},
\end{equation}
and
\begin{equation}\label{iw:F-substitution}
 F_\lambda(\phi(z))
 =\frac{D_\lambda(z)}{(z-1)^2(z+1)^{2k-2}}.
\end{equation}
Also,
\begin{align}
 A_j(0)&=(-1)^{j-1},&
 A_j(1/4)&=j(-1/2)^{j-1},\label{iw:A-endpoints}\\
 F_\lambda(0)&=1-\lambda,&
 F_\lambda(1/4)&=
 \frac{k^2-\lambda(k^2-m^2)}{4^{k-1}}.\label{iw:F-endpoints}
\end{align}
Every root of \(A_j\) is simple, and its root set is
\begin{equation}\label{iw:A-roots}
 \left\{\phi(\zeta):\zeta^j=1,\ \zeta\ne1,-1\right\}.
\end{equation}
Reciprocal values of \(\zeta\) give the same root and are the only
duplication. In particular,
\begin{equation}\label{iw:degree-A}
 \deg A_j=\left\lfloor\frac{j-1}{2}\right\rfloor.
\end{equation}
\end{lemma}

\begin{proof}
At \(t=\phi(z)\), the two roots of \(U^2+U+t\) are
\(\alpha=-z/(1+z)\) and \(\beta=-1/(1+z)\). For \(z\ne\pm1\)
they are distinct. The sequence
\((\alpha^j-\beta^j)/(\alpha-\beta)\) has initial values \(0,1\)
and satisfies \eqref{iw:A}, giving \eqref{iw:A-substitution}.
Substitution in \eqref{iw:F} gives \eqref{iw:F-substitution};
the signs cancel in both products. The endpoint evaluations follow
from \eqref{iw:A} by induction and then from \eqref{iw:F}.

For every finite \(t\ne0,1/4\), the inverse equation
\begin{equation}\label{iw:inverse}
 tz^2+(2t-1)z+t=0
\end{equation}
has two distinct nonzero roots, neither of which is \(1\) or \(-1\).
Its discriminant is \(1-4t\); substitution of \(1\) and \(-1\)
gives \(4t-1\) and \(1\), respectively. At either inverse point,
\begin{equation}\label{iw:phi-derivative}
 \phi'(z)=\frac{1-z}{(1+z)^3}\ne0.
\end{equation}
Since neither endpoint in \eqref{iw:A-endpoints} is zero,
\eqref{iw:A-substitution} now gives every root of \(A_j\).
The numerator \(z^j-1\) has simple roots, the denominator does not
vanish there, and \(\phi\) is unramified there. Thus each root of
\(A_j\) is simple. Equation \eqref{iw:inverse} identifies its two
inverse points as \(z,z^{-1}\). Counting the \(j\)-th roots of unity
after excluding \(1\) and, for even \(j\), \(-1\), proves
\eqref{iw:degree-A}.
\end{proof}

The exceptional points of \(\phi:\mathbb P^1_z\to\mathbb P^1_t\)
will be retained throughout the proof. They satisfy
\[
 \phi(0)=\phi(\infty)=0,\qquad
 \phi(1)=1/4,\qquad \phi(-1)=\infty.
\]
The map is ramified at \(z=1,-1\). More precisely,
\begin{equation}\label{iw:ramification}
 \phi(z)-\frac14=-\frac{(z-1)^2}{4(z+1)^2},
 \qquad
 \operatorname{ord}_{z=-1}\phi=-2.
\end{equation}
For \(\lambda\ne1\), \(D_\lambda(0)=1-\lambda\ne0\), its leading
coefficient is also \(1-\lambda\), and neither inverse point above
\(t=0\) produces a root. The behavior at \(z=1,-1\) is determined
below, rather than inferred from an unramified substitution.

\subsection{Descent to a rational trace}

Set \(w=z^d\), and define
\begin{equation}\label{iw:h}
 h_\lambda(w)
 =(1-\lambda)(w^K+w^{-K})
     +\lambda(w^M+w^{-M})-2,
 \qquad \Theta=w\frac{d}{dw}.
\end{equation}
Then \(H_\lambda(z)=h_\lambda(z^d)\). At every \(z\ne0\), the map
\(z\mapsto z^d\) is unramified, so these two functions have equal
orders of vanishing at corresponding points.

\begin{lemma}[Descent and the tangent equation]\label{iw:descent}
Let \(\lambda\in\mathbb Q\setminus\{0,1\}\). Suppose that
\(w\in\mathbb C^*\), \(w\ne1\), and
\begin{equation}\label{iw:multiple-h}
 h_\lambda(w)=\Theta h_\lambda(w)=0.
\end{equation}
Put \(X=w^K+w^{-K}\) and \(Y=w^M+w^{-M}\). Then \(w\ne-1\),
\(w^K\ne1\), \(w^M\ne1\), and
\begin{equation}\label{iw:branch-values}
 Y=\frac{4K^2}{\lambda\Delta}-2,\qquad
 X=\frac{4M^2}{(\lambda-1)\Delta}-2.
\end{equation}
In particular, \(X,Y\in\mathbb Q\), and
\begin{equation}\label{iw:rational-trace}
 s=w+w^{-1}\in\mathbb Q.
\end{equation}
Furthermore,
\begin{equation}\label{iw:L}
 L_{K,M}(w):=
 (K-M)(w^{K+M}-1)-(K+M)(w^K-w^M)=0.
\end{equation}
\end{lemma}

\begin{proof}
The equations in \eqref{iw:multiple-h} say
\begin{align}
 (1-\lambda)X+\lambda Y&=2,\label{iw:value-eq}\\
 (1-\lambda)K(w^K-w^{-K})
   +\lambda M(w^M-w^{-M})&=0.\label{iw:derivative-eq}
\end{align}
Squaring the second equation and substituting the first gives
\[
 (Y-2)\bigl(\lambda\Delta(Y+2)-4K^2\bigr)=0.
 \tag{\(\ast\)}
\]
If \(Y=2\), then \(w^M=1\); since \(\lambda\ne1\), the first
equation gives \(X=2\), hence \(w^K=1\). Coprimality would imply
\(w=1\), contrary to the hypothesis. This proves the formula for
\(Y\), and substitution gives the formula for \(X\).
The same value equation shows that \(w^K=1\) would force
\(w^M=1\), because \(\lambda\ne0\). Thus neither power equals one.

Because \(K,M\) are coprime, they are not both even. Directly,
\[
 h_\lambda(-1)=
 \begin{cases}
  -4,&K,M\text{ both odd},\\
  -4\lambda,&K\text{ even},\ M\text{ odd},\\
  4(\lambda-1),&K\text{ odd},\ M\text{ even}.
 \end{cases}
\]
This excludes \(w=-1\).

We next prove rationality of the basic trace, without any
assumption that it is nonintegral. Set
\(\alpha=w^K\), \(\beta=w^M\). Since
\(\beta^2-Y\beta+1=0\), the field \(\mathbb Q(\beta)\) has degree
at most two. Equation \eqref{iw:derivative-eq}, in its original
unsquared form, gives
\begin{equation}\label{iw:quadratic-field}
 \alpha=\frac12\left(
 X+\frac{\lambda M}{(\lambda-1)K}
       (\beta-\beta^{-1})\right)\in\mathbb Q(\beta).
\end{equation}
Choose integers \(r,u\) with \(rK+uM=1\). The numbers \(\alpha,\beta\)
are nonzero, so negative powers cause no problem, and
\(w=\alpha^r\beta^u\in\mathbb Q(\beta)\).
If \(\mathbb Q(\beta)=\mathbb Q\), this already proves the claim.
Otherwise its nontrivial conjugation sends \(\beta\) to
\(\beta^{-1}\). By \eqref{iw:quadratic-field}, it sends \(\alpha\)
to \(X-\alpha=\alpha^{-1}\), and therefore sends \(w\) to \(w^{-1}\).
The number \(s=w+w^{-1}\) is fixed by this conjugation and is rational.

Finally, \eqref{iw:value-eq} is equivalent to
\((\lambda-1)(X-2)=\lambda(Y-2)\).
Together with \eqref{iw:derivative-eq}, and with \(X,Y\ne2\),
this gives at the point under consideration
\[
 \frac{\Theta X}{X-2}=\frac{\Theta Y}{Y-2}.
\]
This is a consequence of both equations at the point; no identity
valid only at a point is being differentiated. For any \(j\ge1\)
with \(w^j\ne1\),
\[
 \frac{\Theta(w^j+w^{-j})}{w^j+w^{-j}-2}
 =j\frac{w^j+1}{w^j-1}.
\]
Consequently
\[
 K(w^K+1)(w^M-1)=M(w^M+1)(w^K-1).
\]
Expansion gives \eqref{iw:L}.
\end{proof}

\subsection{Two incompatible denominator divisibilities}

Define \(C_j\in\mathbb Z[S]\) by
\begin{equation}\label{iw:C}
 C_0(S)=2,\qquad C_1(S)=S,\qquad
 C_{j+1}(S)=SC_j(S)-C_{j-1}(S).
\end{equation}
Induction gives
\[
 C_j(w+w^{-1})=w^j+w^{-j},
\]
and \(C_j\) is monic of degree \(j\) for \(j\ge1\).
If \(s=p/q\) is in lowest terms, with \(q>0\), then
\begin{equation}\label{iw:monic-denominator}
 \operatorname{den}(C_j(s))=q^j.
\end{equation}
Indeed, after multiplication by \(q^j\), the numerator is
congruent to \(p^j\) modulo each prime dividing \(q\), and therefore
no factor of \(q\) cancels. This statement includes \(q=1\).

\begin{lemma}[Denominator bound from the tangent equation]
\label{iw:tangent-denominator}
Suppose that \(K>M\ge1\) are coprime, \(w\in\mathbb C^*\),
\(L_{K,M}(w)=0\), and \(w+w^{-1}=p/q\) in lowest terms, with
\(q>0\). Set
\[
 \varepsilon=\gcd(K-M,K+M)\in\{1,2\}.
\]
Then
\begin{equation}\label{iw:upper-divisibility}
 q^M\mid\frac{K-M}{\varepsilon}.
\end{equation}
\end{lemma}

\begin{proof}
There is nothing to prove if \(q=1\). Suppose that a prime
\(\ell\) divides \(q\), and set \(e=v_\ell(q)>0\).
The number \(w\) is algebraic, since it solves
\(W^2-(p/q)W+1=0\). Embed its number field in an algebraic closure
of \(\mathbb Q_\ell\), with \(v_\ell(\ell)=1\).
The equation \(w+w^{-1}=p/q\), whose right side has valuation
\(-e\), implies \(v_\ell(w)\in\{-e,e\}\): if \(v_\ell(w)=0\)
the sum is integral, while otherwise its two summands have
different valuations and the smaller one is the valuation of the sum.
The identity
\[
 L_{K,M}(w^{-1})=-w^{-(K+M)}L_{K,M}(w)
\]
allows replacement by the inverse if necessary, so assume
\(v_\ell(w)=-e\).

Write \(r=K-M\) and \(u=K+M\).
The four nonzero terms
\[
 rw^{K+M},\quad -uw^K,\quad uw^M,\quad -r
\]
have valuations
\[
 v_\ell(r)-(K+M)e,\quad v_\ell(u)-Ke,\quad
 v_\ell(u)-Me,\quad v_\ell(r).
\]
The second is strictly smaller than the third, and the first
is strictly smaller than the fourth. A vanishing sum cannot
have a unique term of smallest valuation. Thus the first two
valuations are equal, giving
\[
 Me=v_\ell(r)-v_\ell(u)>0.
\]
Since \(v_\ell(\varepsilon)=
\min\{v_\ell(r),v_\ell(u)\}=v_\ell(u)\), it follows that
\(v_\ell(q^M)=v_\ell(r/\varepsilon)\).
Combining these conclusions for all \(\ell\mid q\) proves
\eqref{iw:upper-divisibility}.
Finally, \(\varepsilon\) divides \(2\), because
\(\gcd(K-M,M)=1\); it equals two precisely when \(K,M\) are both odd.
\end{proof}

\begin{theorem}[Restriction imposed by the coefficient denominator]
\label{iw:denominator-obstruction}
Let \(\lambda=a/b\in\mathbb Q\setminus\{0,1\}\), with
\(\gcd(a,b)=1\) and \(b>0\). If \(h_\lambda\) has a multiple root
other than \(w=1\), then
\begin{equation}\label{iw:obstruction}
 \boxed{\frac{K+M}{\gcd(K+M,2)}\mid b.}
\end{equation}
In particular, for an integral weight \(\lambda\ne0,1\),
\(h_\lambda\) has no multiple root other than \(1\).
The assertion and its proof remain valid when the rational trace
\(w+w^{-1}\) is an integer.
\end{theorem}

\begin{proof}
Apply Lemma~\ref{iw:descent}, and write \(s=w+w^{-1}=p/q\) in
lowest terms. By \eqref{iw:branch-values},
\[
 C_M(s)+2=\frac{4bK^2}{a\Delta}.
\]
Adding an integer does not change a reduced denominator.
For a prime \(\ell\mid\Delta\), coprimality of \(K\) and \(M\)
implies \(\ell\nmid K\). The \(\ell\)-part of the reduced
denominator of the right side has exponent
\[
 \max\{0,v_\ell(|a|)+v_\ell(\Delta)-v_\ell(4b)\}
 \ \ge\ \max\{0,v_\ell(\Delta)-v_\ell(4b)\}.
\]
Using \eqref{iw:monic-denominator}, this proves
\begin{equation}\label{iw:lower-divisibility}
 \frac{\Delta}{\gcd(\Delta,4b)}\mid q^M.
\end{equation}
On the other hand, Lemma~\ref{iw:tangent-denominator} gives
\eqref{iw:upper-divisibility}. Write
\(g_0=\gcd(\Delta,4b)\), \(r=K-M\), \(u=K+M\). Since \(\Delta=ru\),
the two divisibilities imply
\[
 \frac{ru}{g_0}\mid\frac r\varepsilon,
 \qquad\text{hence}\qquad
 \varepsilon u\mid g_0\mid4b.
\]
If \(K,M\) have opposite parity, then \(\varepsilon=1\) and
\(u\) is odd, so \(u\mid b\). If they are both odd, then
\(\varepsilon=2\), \(u\) is even, and \(u/2\mid b\).
These are exactly \eqref{iw:obstruction}.
Neither argument requires \(q>1\): for \(q=1\) the upper
divisibility is immediate and the lower one is unchanged.

For \(b=1\), the obstruction is impossible. In the odd case
\(u\ge3\), and in the even case \(u\ge4\), so the positive
integer \(u/\gcd(u,2)\) is always at least two.
\end{proof}

\subsection{The resonant coefficient and exact multiplicities}

\begin{proposition}[All possible multiplicities]\label{iw:multiplicities}
Let \(\lambda\in\mathbb Q\setminus\{0,1\}\), and put
\begin{equation}\label{iw:resonance}
 \lambda_*=\frac{K^2}{K^2-M^2}.
\end{equation}
The root \(w=1\) of \(h_\lambda\) has multiplicity two when
\(\lambda\ne\lambda_*\), and multiplicity four when
\(\lambda=\lambda_*\). Every multiple root other than \(1\)
has multiplicity exactly two. Such roots do not exist at
\(\lambda=\lambda_*\).

For the original polynomial, \(t=0\) is never a root;
\(t=1/4\) is a root exactly when \(\lambda=\lambda_*\), and then
it is simple. In the resonant case,
\begin{equation}\label{iw:resonant-gcd}
 \gcd(F_{\lambda_*},F_{\lambda_*}')\doteq A_d^3,
\end{equation}
and \(F_{\lambda_*}/A_d^4\) is squarefree and coprime to \(A_d\).
The number \(\lambda_*\) is not an integer.
\end{proposition}

\begin{proof}
At \(w=1\),
\[
 h_\lambda=\Theta h_\lambda=0,\qquad
 \Theta^2h_\lambda=2(K^2-\lambda\Delta).
 \tag{\(\dagger\)}
\]
The local coordinate \(w=e^u\) identifies \(\Theta\) with
\(d/du\), and
\[
 h_\lambda(e^u)
 =2\bigl((1-\lambda)\cosh(Ku)+\lambda\cosh(Mu)-1\bigr)
\]
is even in \(u\). At \(\lambda=\lambda_*\), its second
derivative vanishes, while its fourth derivative is
\[
 2\bigl((1-\lambda_*)K^4+\lambda_*M^4\bigr)
 =-2K^2M^2\ne0.
 \tag{\(\ddagger\)}
\]
This proves the claimed multiplicities at \(1\).

At a multiple root different from \(1\), substitute
\eqref{iw:branch-values} into
\[
 \Theta^2h_\lambda=(1-\lambda)K^2X+\lambda M^2Y
\]
to obtain
\begin{equation}\label{iw:extra-second-derivative}
 \Theta^2h_\lambda=2(\lambda\Delta-K^2).
\end{equation}
If \(\lambda=\lambda_*\), the formula for \(Y\) in
\eqref{iw:branch-values} gives \(Y=2\), which was excluded in
Lemma~\ref{iw:descent}. Otherwise
\eqref{iw:extra-second-derivative} is nonzero. At a multiple
root \(w\ne0\), one has \(\Theta^2h=w^2h''\), so that root is
exactly double.

Equation \eqref{iw:F-endpoints} handles \(t=0\), and shows that
\(1/4\) can be a root only at \(\lambda_*\). At this coefficient
\(D_\lambda\) has order four at \(z=1\): multiplication by \(z^k\)
and the map \(z\mapsto z^d\) do not change the order.
The denominator in \eqref{iw:F-substitution} has order two there,
so \(F_{\lambda_*}(\phi(z))\) has order two at \(z=1\).
By \eqref{iw:ramification}, \(F_{\lambda_*}\) has a simple root
at \(t=1/4\).

At every other finite root, the inverse points from
\eqref{iw:inverse} lie outside \(0,1,-1\), and \(\phi\) is
unramified. The multiplicity of \(F_\lambda\) therefore equals
that of \(D_\lambda\), and hence that of \(h_\lambda\).
At \(\lambda_*\), all multiple roots of \(D_\lambda\) satisfy
\(z^d=1\), and have multiplicity four. Their finite images other
than \(1/4\) are exactly the simple roots of \(A_d\), by
Lemma~\ref{iw:substitution}. The point \(z=-1\) is over infinity
and contributes no finite root. This proves
\eqref{iw:resonant-gcd} and the assertions about the quotient.

Finally, \(\gcd(K^2,\Delta)=1\) and \(\Delta\ge3\), so the
reduced denominator of \(\lambda_*\) is \(\Delta>1\).
\end{proof}

\begin{corollary}[Restriction for the original polynomial]
\label{iw:restriction-F}
If \(\lambda=a/b\ne0,1\) is in lowest terms and \(F_\lambda\)
has a repeated root that is not a root of \(A_d\), then
\eqref{iw:obstruction} holds.
\end{corollary}

\begin{proof}
The points \(0\) and \(1/4\) cannot be repeated roots by
Proposition~\ref{iw:multiplicities}. Every remaining repeated
root lifts unramified to a multiple root of \(D_\lambda\).
If its image were \(w=z^d=1\), it would be a root of \(A_d\).
Thus Theorem~\ref{iw:denominator-obstruction} applies.
\end{proof}

\subsection{An exact description for rational weights}

The preceding results also give a finite description of all
additional repeated factors, without asserting that the obstruction
\eqref{iw:obstruction} alone is sufficient.
Define
\begin{equation}\label{iw:R}
 R_{K,M,\lambda}(S)
 =(1-\lambda)C_K(S)+\lambda C_M(S)-2,
\end{equation}
and, for \(\lambda\ne0,1,\lambda_*\), set
\begin{equation}\label{iw:S}
 \mathcal S_{K,M,\lambda}
 =\{s\in\mathbb Q\setminus\{-2,2\}:
         R_{K,M,\lambda}(s)=R_{K,M,\lambda}'(s)=0\}.
\end{equation}
For \(s\ne\pm2\), put
\begin{equation}\label{iw:E}
 E_{d,s}(t)=t^d\bigl(C_d(t^{-1}-2)-s\bigr).
\end{equation}

\begin{theorem}[Rational-weight repeated factors]\label{iw:rational}
For \(\lambda\in\mathbb Q\setminus\{0,1,\lambda_*\}\),
\(\mathcal S=\mathcal S_{K,M,\lambda}\) is finite, and
\begin{equation}\label{iw:rational-gcd}
 \gcd(F_\lambda,F_\lambda')
 \doteq A_d\prod_{s\in\mathcal S}E_{d,s}.
\end{equation}
Each factor \(E_{d,s}\) is a squarefree polynomial in
\(\mathbb Q[t]\) of degree \(d\), with constant term \(1\).
These factors are pairwise coprime and are coprime to \(A_d\).
All roots represented in \eqref{iw:rational-gcd} are double
roots of \(F_\lambda\), and all other roots of \(F_\lambda\)
are simple.

If \(\mathcal S\ne\varnothing\) and \(\lambda=a/b\) is reduced,
then \eqref{iw:obstruction} holds. Hence, for a fixed rational
weight, only finitely many primitive pairs \(K>M\) can have
additional repeated factors; the common multiplier \(d\) remains
unrestricted.
\end{theorem}

\begin{proof}
The identity \(h_\lambda(w)=R_{K,M,\lambda}(w+w^{-1})\) gives
\[
 \Theta h_\lambda(w)
 =(w-w^{-1})R_{K,M,\lambda}'(w+w^{-1}).
 \tag{\(\star\)}
\]
A multiple root of \(h_\lambda\) other than \(1\) is not \(-1\),
has a rational trace by Lemma~\ref{iw:descent}, and therefore
determines an element of \(\mathcal S\). Conversely, for
\(s\in\mathcal S\), both roots of \(w^2-sw+1=0\) are different
from \(0,1,-1\), and \((\star)\) makes them multiple roots of
\(h_\lambda\). Proposition~\ref{iw:multiplicities} shows that
they are exactly double. The set is finite because
\(R_{K,M,\lambda}\) has degree \(K\), with nonzero leading
coefficient \(1-\lambda\). In particular, this argument also
proves that every multiple root of \(R_{K,M,\lambda}\) other
than \(2\) is rational; rationality is not an additional
restriction that discards possible repeated factors.

The expression \eqref{iw:E} is a polynomial because \(C_d\)
has degree \(d\). Its constant term is \(1\), and its coefficient
of \(t^d\) is \(C_d(-2)-s=2(-1)^d-s\ne0\).
For \(t=\phi(z)\ne0\),
\[
 t^{-1}-2=z+z^{-1},\qquad
 C_d(t^{-1}-2)=z^d+z^{-d}.
\]
Thus the roots of \(E_{d,s}\) are the images of the \(2d\)
distinct solutions of \(z^d+z^{-d}=s\). None equals \(1\) or
\(-1\), and reciprocal pairs give exactly \(d\) distinct
finite images. The logarithmic derivative of this trace is
\(d(z^d-z^{-d})\ne0\), and \(\phi'\ne0\) there, proving
squarefreeness. Different \(s\)'s have disjoint root sets.
A root of \(A_d\) has trace \(z^d+z^{-d}=2\), so it is not
a root of any \(E_{d,s}\).

The endpoint and multiplicity analysis in
Proposition~\ref{iw:multiplicities} now proves
\eqref{iw:rational-gcd}, including the exact multiplicities.
The denominator condition follows from
Theorem~\ref{iw:denominator-obstruction}. If \(K+M\) is odd,
it divides \(b\); if it is even, it divides \(2b\).
There are finitely many such positive sums and hence finitely
many primitive pairs. Formula \eqref{iw:E} treats every
positive common multiplier \(d\).
\end{proof}

\subsection{Degrees, normalization, and completion of the integral theorem}

\begin{proposition}[Degrees and geometric genus]\label{iw:geometry}
For \(\lambda\in\mathbb Q\setminus\{0,1\}\),
\begin{equation}\label{iw:degree-F}
 \deg F_\lambda=
 \begin{cases}
 k-1,&d\text{ odd},\\
 k-2,&d\text{ even and }\lambda\ne\lambda_*,\\
 k-3,&d\text{ even and }\lambda=\lambda_*.
 \end{cases}
\end{equation}
In the resonant case,
\[
 Q_\lambda=F_{\lambda_*}/A_d^4
\]
is squarefree, has degree \(k-2d+1\ge1\), and the smooth
projective normalization of \(y^2=F_{\lambda_*}(t)\) has genus
\begin{equation}\label{iw:resonant-genus}
 \left\lfloor\frac{k-2d}{2}\right\rfloor.
\end{equation}

For \(\lambda\ne0,1,\lambda_*\), let
\(\nu=|\mathcal S_{K,M,\lambda}|\). The polynomial
\begin{equation}\label{iw:squarefree-model}
 Q_\lambda=
 \frac{F_\lambda}
 {A_d^2\prod_{s\in\mathcal S_{K,M,\lambda}}E_{d,s}^2}
\end{equation}
is squarefree, and
\begin{equation}\label{iw:general-branch-degree}
 N:=\deg Q_\lambda=k-d-2d\nu\ge0.
\end{equation}
If \(N\ge1\), the curve is geometrically integral and its
smooth projective normalization has genus
\(\lfloor(N-1)/2\rfloor\).
If \(N=0\), its normalization over \(\mathbb C\) consists
of two copies of \(\mathbb P^1\); a genus formula for a single
geometrically integral double cover is not applicable.
\end{proposition}

\begin{proof}
All possible degree losses are visible at \(z=-1\).
The denominator in \eqref{iw:F-substitution} has order \(2k-2\)
there, and \(\phi\) has a pole of order two.
If \(d\) is odd, \(w=(-1)^d=-1\) and
Lemma~\ref{iw:descent}'s direct evaluation of \(h_\lambda(-1)\)
is nonzero; hence \(D_\lambda(-1)\ne0\).
If \(d\) is even, \(w=1\), and
Proposition~\ref{iw:multiplicities} gives order two for
\(D_\lambda\), or order four at \(\lambda_*\).
For any nonzero polynomial \(F\) of degree \(n\),
\(\operatorname{ord}_{z=-1}F(\phi(z))=-2n\).
Consequently, if the order of \(D_\lambda\) at \(-1\) is
\(j\in\{0,2,4\}\), then
\[
 -2\deg F_\lambda=j-(2k-2).
 \tag{\(\diamond\)}
\]
This proves \eqref{iw:degree-F} and accounts explicitly for
the pole over \(t=\infty\).

By \eqref{iw:degree-A}, \(2\deg A_d\) is \(d-1\) for odd \(d\),
and \(d-2\) for even \(d\). Using \eqref{iw:degree-F} gives
\(\deg(F_{\lambda_*}/A_d^4)=k-2d+1\) in both parities.
Since \(k=dK\) with \(K\ge2\), this degree is at least one.
In the nonresonant case, Theorem~\ref{iw:rational} and
\(\deg E_{d,s}=d\) give \eqref{iw:general-branch-degree}.
The quotient is nonzero, so its degree is nonnegative.

In either case, removal of the displayed square factors is
a birational transformation to \(v^2=Q_\lambda(t)\).
If \(Q_\lambda\) is squarefree of degree \(N\ge1\), it is not
a square in \(\mathbb C(t)\), so the resulting curve is
geometrically integral. Its double cover of the projective
\(t\)-line is ramified at the \(N\) simple finite roots, and
at infinity exactly when \(N\) is odd. There are no other
branch points. Riemann--Hurwitz \cite[Tag~0C1B]{Stacks} gives
\[
 2g-2=-4+
 \begin{cases}
 N,&N\text{ even},\\
 N+1,&N\text{ odd},
 \end{cases}
\]
so \(g=\lfloor(N-1)/2\rfloor\).
For the resonant degree \(N=k-2d+1\), this is
\eqref{iw:resonant-genus}.
If \(N=0\), then \(Q_\lambda=c\in\mathbb Q^*\) and, over
\(\mathbb C\), the equation is \(v=\pm\sqrt c\).
Each component is a rational curve, with projective
normalization \(\mathbb P^1\).
\end{proof}

\begin{proof}[Proof of Theorem~\ref{iw:integral}]
For an integral \(\lambda\ne0,1\),
Theorem~\ref{iw:denominator-obstruction} excludes every
multiple root of \(h_\lambda\) except \(1\).
Proposition~\ref{iw:multiplicities} shows that the resonant
coefficient is not integral. Thus the roots of \(D_\lambda\)
at \(z^d=1\) are exactly double. The point \(t=0\) is not a
root, and \(t=1/4\) is not a root because the coefficient
is nonresonant. The unramified inverse points
\eqref{iw:inverse} transfer every other finite multiplicity
exactly. Lemma~\ref{iw:substitution} identifies the finite
repeated roots with the simple roots of \(A_d\), proving
\eqref{iw:integral-gcd}. The point \(z=-1\) lies over
infinity and has already been included in the degree
calculation; \(z=0,\infty\) lie over the nonroot \(t=0\).
No finite root is omitted.

The root classification gives \(A_d^2\mid F_\lambda\)
over \(\mathbb C[t]\), hence over \(\mathbb Q[t]\) by
polynomial division. The quotient is squarefree and
coprime to \(A_d\). Because \(A_d(0)=\pm1\), \(A_d\) is
primitive in \(\mathbb Z[t]\). Gauss's lemma therefore
shows that the quotient belongs to \(\mathbb Z[t]\).
In Theorem~\ref{iw:rational} the set \(\mathcal S\) is
empty, so Proposition~\ref{iw:geometry} gives degree
\(k-d>0\) and the genus in \eqref{iw:integral-degree-genus}.
\end{proof}

\subsection{The two degenerate weights and rational examples}

\begin{proposition}[Weights zero and one]\label{iw:degenerate}
For \(k>m\ge1\),
\begin{align}
 F_0&=A_k^2,&
 \gcd(F_0,F_0')&\doteq A_k,\label{iw:lambda-zero}\\
 F_1&=t^{k-m}A_m^2,&
 \gcd(F_1,F_1')&\doteq t^{k-m-1}A_m.\label{iw:lambda-one}
\end{align}
These formulas include constant \(A_k\) or \(A_m\).
The degrees are \(2\lfloor(k-1)/2\rfloor\) and
\(k-m+2\lfloor(m-1)/2\rfloor\), respectively.
Over \(\mathbb C\), the curve for \(\lambda=0\) has two
rational components. For \(\lambda=1\), it has two rational
components when \(k-m\) is even, and is geometrically
integral and rational when \(k-m\) is odd.
\end{proposition}

\begin{proof}
The first identity is the definition. For the second,
\eqref{iw:D} at \(\lambda=1\) simplifies to
\[
 D_1(z)=z^{k-m}(z^m-1)^2.
 \tag{\(\triangle\)}
\]
Together with \eqref{iw:A-substitution} and
\eqref{iw:F-substitution}, this gives
\(F_1(\phi(z))=\phi(z)^{k-m}A_m(\phi(z))^2\).
It is therefore an identity of polynomials in \(t\).
The roots of \(A_j\) are simple and \(A_j(0)\ne0\), so
the stated gcds follow by their exact root multiplicities.
The degree formulas follow from \eqref{iw:degree-A}.

For \(\lambda=0\), the equation factors as
\((y-A_k)(y+A_k)=0\).
For \(\lambda=1\), remove \(A_m^2\) and the largest even
power of \(t\). The remaining equation is either \(v^2=1\)
or \(v^2=t\), according as \(k-m\) is even or odd.
The asserted normalizations follow.
\end{proof}

The denominator obstruction is necessary, not sufficient,
and additional repeated roots do occur for rational weights.
For example, direct substitution in \eqref{iw:A} gives
\begin{align}
 F_{4/5;3,2}(t)&=\frac{(t+1)^2}{5},\label{iw:example-one}\\
 F_{32/5;4,1}(t)&=\frac{(2t-3)^2(8t-3)}5,\label{iw:example-two}\\
 F_{25/9;5,1}(t)&=
   \frac{(3t-2)^2(t^2+12t-4)}9.\label{iw:example-three}
\end{align}
All three pairs have \(d=1\). Their additional rational traces
are respectively \(-3,-4/3,-1/2\).
The last coefficient is the rational square \((5/3)^2\);
thus the integral-square consequence cannot be extended
without qualification to every rational square.
As a resonant example,
\[
 F_{4/3;6,3}(t)=\frac{(t-1)^4(4t-1)}3,
\]
which displays both the fourth-order common root and the
simple root at \(t=1/4\).

\begin{remark}[Extent of the generalization]\label{iw:scope}
The integral-weight theorem is a statement for every
integer coefficient and every admissible pair of exponents,
apart from the two completely described degenerate weights.
No numerical bound on the coefficient or the exponents is
used in its proof. For rational weights the stronger
denominator restriction and the exact factor description
are separate assertions: the former narrows the possible
primitive exponent pairs, while the latter determines the
actual repeated factors for each pair. These results do not
by themselves certify the historical novelty of the
weighted statements or resolve a different named conjecture.
\end{remark}


\section{Method transfer and the six research gates}\label{sec:gates}

The previous corpus calculations used exact valuation and
ramification information, including the Eisenstein shift
$\Phi_{2^h}(1+X)$, to determine a field requirement in a paired
cancellation system. The preceding quadrinomial manuscript
\cite{Previous} used rational traces and exact reduced denominators
to exclude repeated roots. The present two applications retain
those verified mechanisms in distinct settings:
\begin{enumerate}
 \item In the cyclotomic trace family, translation gives an
 Eisenstein derivative. The derivative equation then isolates a
 unique least-valuation term in a critical value. Its denominator
 forces the full algebraic degree, or equivalently makes the
 transformed discriminant Eisenstein.
 \item In the weighted recurrence family, two coprime traces and
 the unsquared derivative equation give a rational fundamental
 trace. A rational coefficient denominator must simultaneously
 satisfy two opposing divisibility conditions. For integer
 weights the conditions are incompatible.
\end{enumerate}
The rational reconstruction formula of Section~\ref{sec:inverse}
also supplies a second, global algebraic route in the first
application. None of these transfers assumes a solution of
Kaplansky's zero-divisor conjecture. No result about arbitrary
torsion-free groups is asserted here.

\Needspace{27\baselineskip}
\subsection{G1--G6 audit}
\renewcommand{\arraystretch}{1.15}
\begin{longtable}{@{}p{0.07\linewidth}p{0.20\linewidth}p{0.66\linewidth}@{}}
\toprule
Gate & Assessment & Evidence and limit\\
\midrule
G1 & Existing named target & Jones's Conjecture~3.3 was published in
2021, with the exact family and normalized factors stated on p.~1476
of \cite{Jones}.\\
G2 & Statement retained & Theorem~\ref{thm:jones} uses the original
polynomials, normalization, and prime range, with the integer
parameter convention explicitly used in \cite{BNW}. Every change
of variable used for irreducibility is invertible over the
specified field.\\
G3 & Primary-source evidence & The 2026 published follow-up
\cite{BNW} states a partial proof and proves the factorization
formula. We did not identify an earlier complete proof of the
remaining assertions in the sources examined. This is a bounded
literature finding, not an exhaustive historical theorem.\\
G4 & Full proof & Theorem~\ref{thm:criterion} and the proof of
Theorem~\ref{thm:jones} cover every odd prime, both reciprocal families, and every nonzero
integer $r$, including $q=3$. There is no unproved conjectural
assumption or finite computation in the proof.\\
G5 & Substantive range & The exact critical-value valuation establishes
irreducibility for unbounded degree, beyond the already proved
factorization. The integer-weight theorem also extends the earlier
weight-$4$ result to all nondegenerate integer weights. The basic
valuation and Eisenstein tools are classical.\\
G6 & Not certified & No identical full solution was found in the
focused search, but negative search results cannot certify
priority. External literature review and mathematical review
remain necessary for a historical-first claim.\\
\bottomrule
\end{longtable}
\renewcommand{\arraystretch}{1}

\clearpage
\appendix
\section{Exact computational checks}\label{sec:checks}

The all-prime proof uses no computation. As checks on signs,
normalization, and boundary cases, exact symbolic arithmetic was
performed with SymPy 1.14.0. For
\[
 q\in\{5,7,11,13,17,19,23,29,31\},
\]
the trace identity, shifted congruence, exact first coefficient,
degree and constant of $D_q$, all Eisenstein coefficient conditions
for $E_q$, and $E_q(1)=1$ were verified. For
$q\in\{3,5,7,11\}$ and both $a$, the discriminant of the original
reciprocal polynomial was computed with $r,t$ symbolic and
compared with \eqref{eq:jones-zero}--\eqref{eq:jones-one}.
All checks passed.

For the integer-weight theorem, exact gcd calculations were checked
for all $2\le k\le15$, $1\le m<k$, and
$\lambda\in\{-5,-4,-3,-2,-1,2,3,\ldots,10\}$, a total of
$1470$ cases. All passed. These checks test finite examples;
the proofs in the main text establish the unrestricted statements.

The following compact verification program is included so that the
single \texttt{.tex} file contains a reproducible check as well as
the mathematical argument. It requires Python and SymPy; it is not
required to compile the document.

\begin{lstlisting}[language=Python]
import sympy as S
u, Z, Y, t, r = S.symbols("u Z Y t r")
C = [S.Integer(2), u]
for j in range(1, 15):
    C.append(S.expand(u*C[-1] - C[-2]))
for q in [5, 7, 11, 13, 17, 19, 23, 29, 31]:
    d = (q-1)//2
    H = 1 + sum(C[1:d+1])
    P = S.Poly(S.expand(H.subs(u, u+2)), u)
    assert all(c % q == 0 for c in P.all_coeffs()[1:])
    assert S.diff(H, u).subs(u, 2) == q*(q*q-1)//24
    D = S.discriminant(H+Z, u)
    E = S.Poly(S.expand(D.subs(Z, q*(Y-1))/q**(d-1)),
               Y, domain=S.ZZ)
    assert E.degree() == d-1 and E.LC() % q != 0
    assert all(c % q == 0 for c in E.all_coeffs()[1:])
    assert E.TC() % (q*q) != 0 and E.eval(1) == 1
    h = S.Poly(E.as_expr().subs(Y, 1+4*r*q*t), t, r)
    assert h.content() == 1
A = [S.Integer(0), S.Integer(1)]
for j in range(1, 30):
    A.append(S.expand(-A[-1]-t*A[-2]))
count = 0
for k in range(2, 16):
    for m in range(1, k):
        for lam in list(range(-5, 0)) + list(range(2, 11)):
            F = S.Poly(A[k]**2-lam*A[k-m]*A[k+m], t,
                       domain=S.QQ)
            G = S.gcd(F, F.diff()).monic()
            expected = S.Poly(A[S.igcd(k, m)], t,
                              domain=S.QQ).monic()
            assert G == expected
            count += 1
assert count == 1470
print("Exact finite checks passed; see text for global proofs.")
\end{lstlisting}

\clearpage
\section{Literature search and document provenance}\label{sec:provenance}

The focused search on 23 November 2026 examined Jones's original
paper, Barman--Narode--Wagh's published abstract and public full
manuscript, and Narode's related 2026 paper. Searches used the
original title and conjecture number together with
\emph{irreducibility}, \emph{critical values},
\emph{reciprocal discriminant}, and \emph{Eisenstein}.
No identical proof of the normalized factors' all-prime
irreducibility was identified in the sources examined. Other
results described as proofs of a Jones conjecture must be matched
to their actual statement: the conjecture concerning
$x^n+A(Bx+1)^m$ is a different problem.

The exact conjecture is cited from the original 2021 journal
article. The 2026 follow-up's explicit integral-parameter
formulation is also cited, preventing the ambiguous phrase
``nonzero constant'' from being used to claim false integrality
over all rational parameters. The content and sign conventions
are fixed in the text. No claim of priority relies on a lack of
search results.

This manuscript was developed with ChatGPT assistance. Independent
internal proof checks and symbolic checks were used to audit the
argument; they do not constitute external peer review. The PDF
is compiled from this one \texttt{.tex} file, which contains the
entire text, bibliography, and verification program. There are no
external figure assets, bibliography databases, or
\texttt{\string\input} dependencies. Preparation of the files does
not imply an external submission or a public priority claim.

\begin{thebibliography}{9}

\bibitem{Jones}
Lenny Jones,
\emph{Infinite families of reciprocal monogenic polynomials and
their Galois groups},
New York Journal of Mathematics \textbf{27} (2021), 1465--1493,
especially Lemma~3.2, Conjecture~3.3, and Theorem~3.4.
\url{https://nyjm.albany.edu/j/2021/27-57p.pdf}.

\bibitem{BNW}
Rupam Barman, Anuj Narode, and Vinay Wagh,
\emph{On monogeneity of reciprocal polynomials},
The Ramanujan Journal \textbf{69} (2026), Article~56,
published 12 February 2026.
DOI: \href{https://doi.org/10.1007/s11139-026-01328-2}{10.1007/s11139-026-01328-2}.
Public manuscript: arXiv:2601.22453v1, especially
equations~(1.4)--(1.6) and Theorem~1.3.
\url{https://arxiv.org/html/2601.22453v1}.

\bibitem{Narode}
Anuj Narode,
\emph{Necessary and sufficient condition for reciprocal polynomials
to be monogenic}, arXiv:2609.33766, 2026.
\url{https://arxiv.org/pdf/2609.33766}.

\bibitem{BU}
Andrew Bremner and Maciej Ulas,
\emph{Some observations concerning reducibility of quadrinomials},
arXiv:1310.5346v1, 20 October 2013.
Journal version: Acta Mathematica Hungarica \textbf{145} (2015),
320--349, DOI: 10.1007/s10474-015-0478-9.
\url{https://arxiv.org/pdf/1310.5346}.

\bibitem{KSL}
Korea Superintelligence Labs (generated by Machina Mathematica),
\emph{$A_{\gcd(m,n)}^2$ always divides $F_{m,n}$: multiple roots of
the Bremner--Ulas quadrinomial discriminants}, Version~1, 2026,
especially Conjecture~A.
\url{https://ideosphere.ai/papers/quadrinomial-sqfree-v1.pdf}.

\bibitem{Previous}
\emph{Exact Multiplicities in the Bremner--Ulas Quadrinomial
Discriminants: A proof of the published corrected gcd conjecture},
research manuscript prepared for Keio, 23 November 2026.
Unpublished preceding manuscript in this research session.
The present integer-weight proof is self-contained and does not
require access to that manuscript.

\bibitem{Stacks}
The Stacks Project Authors,
\emph{The Stacks Project}, Section~53.12, Riemann--Hurwitz,
Tag~0C1B, consulted 23 November 2026.
\url{https://stacks.math.columbia.edu/tag/0C1B}.

\end{thebibliography}
\end{document}
